\subsection{MAGNUS GROUP AND FREE GROUP}

THEOREM 1. Let \(r\) be a mapping of \(\mathbf{X}\) into \(\hat{\mathbf{A}}(\mathbf{X})\) such that \(\omega(r(x)) \geqslant 2\) for all \(x \in \mathbf{X}\) . The unique homomorphism \(g\) of the free group \(\mathrm{F}(\mathbf{X})\) into the Magnus group \(\Gamma(\mathbf{X})\) such that \(g(x) = 1 + x + r(x)\) for all \(x \in \mathbf{X}\) is injective.

We first prove three lemmas.

Lemma 2. Let \(n\) be a non-zero rational integer. In the ring of formal power series \(\mathbf{K}[[t]]\) we write \((1 + t)^n = \sum_{j \geqslant 0} c_{j,n} t^j\) . There exists an integer \(j \geqslant 1\) such that \(c_{j,n} \neq 0\) .

If \(n > 0\) , then \(c_{n,n} = 1\) by the binomial formula.

Suppose that \(n < 0\) and let \(m = -n\) . If \(c_{j,n} = 0\) for all \(j \geqslant 1\) , then \((1 + t)^n = 1\) , whence, taking the inverse, \((1 + t)^m = 1\) , which contradicts the formula \(c_{m,m} = 1\) .

Lemma 3. Let \(x_1, \ldots, x_s\) be elements of \(\mathbf{X}\) such that \(s \geqslant 1\) and \(x_i \neq x_{i+1}\) for \(1 \leqslant i \leqslant s-1\) ; let \(n_1, \ldots, n_s\) be non-zero rational integers. Then the element \(\prod_{i=1}^{s} (1 + x_i)^{n_i}\) of \(\hat{\mathbf{A}}(\mathbf{X})\) is \(\neq 1\) .

Let \(m\) be a maximal ideal of \(K\) and \(k\) the field \(K/m\) ; let \(p: \hat{A}_K(X) \to \hat{A}_k(X)\) be the unique continuous homomorphism of unital \(K\) -algebras such that \(p(x) = x\) for \(x \in X\) (no. 1, Proposition 1). It suffices to prove that \(p\left(\prod (1 + x_i)^{n_i}\right) \neq 1\) and the problem is reduced to the case where \(K\) is a field.

In the notation of Lemma 2:

\[
\prod_ {i = 1} ^ {s} (1 + x _ {i}) ^ {n _ {i}} = \sum_ {b _ {i} \geqslant 0} c _ {b _ {1}, n _ {1}} \dots c _ {b _ {s}, n _ {s}} x _ {1} ^ {b _ {1}} \dots x _ {s} ^ {b _ {s}}.
\]

By Lemma 2, there exist integers \(a_i > 0\) such that \(c_{a_i, n_i} \neq 0 (1 \leqslant i \leqslant s)\) . By Algebra, Chapter I, §7, no. 4, Proposition 6, no monomial \(x_1^{b_1} \ldots x_s^{b_s}\) such that \(b_{i} \geqslant 0\) and \((b_{1}, \ldots, b_{s}) \neq (a_{1}, \ldots, a_{s})\) can be equal to \(x_{1}^{a_{1}} \ldots x_{s}^{a_{s}}\) . It follows that the coefficient of \(x_{1}^{a_{1}} \ldots x_{s}^{a_{s}}\) in \(\prod_{i=1}^{s} (1 + x_{i})^{n_{i}}\) is \(c_{a_{1}, n_{1}} \ldots c_{a_{s}, n_{s}} \neq 0\) , which implies the result.

Lemma 4. Let \(\sigma\) be the continuous endomorphism of \(\hat{\mathbf{A}}(\mathbf{X})\) such that \(\sigma(x) = x + r(x)\) for \(x \in \mathbf{X}\) (no. 1, Proposition 1). Then \(\sigma\) is an automorphism and \(\sigma(\hat{\mathbf{A}}_m(\mathbf{X})) = \hat{\mathbf{A}}_m(\mathbf{X})\) for all \(m \in \mathbf{N}\) .

\(\sigma(x) \equiv x \bmod. \hat{A}_2(X)\) for \(x \in X\) , whence, for \(n \geqslant 1\) and \(x_1, \ldots, x_n\) in \(X\) ,

\[
\sigma (x _ {1}) \dots \sigma (x _ {n}) \equiv x _ {1} \dots x _ {n} \mod \hat {\mathrm{A}} _ {n + 1} (\mathrm{X});
\]

it follows by linearity that \(\sigma(a) \equiv a\) modulo \(\hat{\mathbf{A}}_{n+1}(\mathbf{X})\) for all \(a \in \mathbf{A}^n(\mathbf{X})\) and in particular \(\sigma(\mathbf{A}^n(\mathbf{X})) \subset \hat{\mathbf{A}}_n(\mathbf{X})\) . It follows that \(\sigma(\hat{\mathbf{A}}^m(\mathbf{X})) \subset \hat{\mathbf{A}}_n(\mathbf{X})\) for \(m \geqslant n\) , whence \(\sigma(\hat{\mathbf{A}}_n(\mathbf{X})) \subset \hat{\mathbf{A}}_n(\mathbf{X})\) . In other words, \(\sigma\) is compatible with the filtration \((\hat{\mathbf{A}}_m(\mathbf{X}))\) on \(\mathbf{A}(\mathbf{X})\) and its restriction to the associated graded ring is the identity. Hence \(\sigma\) is bijective (Commutative Algebra, Chapter III, §2, no. 8, Corollary 3 to Theorem 1).

Finally we prove Theorem 1. Let \(w \neq 1\) be an element of \(\mathrm{F}(\mathbf{X})\) . By Algebra, Chapter I, § 7, no. 5, Proposition 7, there exist \(x_{1}, \ldots, x_{s}\) in \(\mathbf{X}\) and non-zero rational integers \(n_{1}, \ldots, n_{s}\) such that \(s \geqslant 1\) , \(x_{i} \neq x_{i+1}\) ( \(1 \leqslant i \leqslant s - 1\) ) and

\[
w = x _ {1} ^ {n _ {1}} \dots x _ {s} ^ {n _ {s}}.
\]

In the notation of Lemma 4,

\[
g (w) = \prod (1 + \sigma (x _ {i})) ^ {n _ {i}} = \sigma (\prod (1 + x _ {i}) ^ {n _ {i}}),
\]

hence \(g(w) \neq 1\) by Lemmas 3 and 4.

